\section{Methodology}
\label{sec:methodology}

This section details the proposed Cost-Sensitive Adaptive Random Forest (CS-ARF) architecture. We first describe the foundational mechanism of the standard Adaptive Random Forest and its concept drift detector. Subsequently, we formulate the mathematical basis of our proposed Implicit Online Oversampling technique to handle extreme class imbalance in a single-pass stream.

\subsection{Standard Adaptive Random Forest and Concept Drift}
The Adaptive Random Forest (ARF) algorithm is an ensemble learning method specifically designed for high-velocity data streams. ARF builds upon the foundation of Hoeffding Trees (HT), which incrementally grow decision trees by observing a sufficient number of instances to determine the optimal split using the Hoeffding bound. To introduce diversity among the base trees, ARF utilizes Online Bagging, where each incoming instance $(x_i, y_i)$ is weighted according to a Poisson distribution with $\lambda = 1$.

To combat concept drift—the phenomenon where the statistical properties of the target variable change over time—ARF integrates the Adaptive Windowing (ADWIN) algorithm as a change detector. ADWIN continuously monitors the prequential error rate of each tree. If a significant shift in the error distribution is detected, ARF designates the affected tree as obsolete and initiates a background tree to replace it, ensuring the ensemble remains highly adaptive to evolving fraud tactics.

\subsection{The Challenge of Extreme Imbalance in Hoeffding Trees}
In credit card fraud detection, the data stream is extremely skewed, with legitimate transactions often outnumbering fraudulent ones by a ratio of $1000:1$ or more. Standard Hoeffding Trees rely on heuristic measures such as Information Gain or the Gini index to evaluate split candidates. 

Let the stream be defined as a sequence of transactions $S = \{ (x_1, y_1), (x_2, y_2), \dots, (x_t, y_t) \}$, where $y_t \in \{0, 1\}$ represents legitimate and fraudulent classes, respectively. When the prior probability of fraud $P(y=1) \approx 0.002$, the entropy reduction achieved by isolating the fraud class is mathematically negligible. Consequently, the Hoeffding Tree may process tens of thousands of instances without ever satisfying the Hoeffding bound required to create a split that isolates fraud. The model effectively collapses into a majority-class predictor, rendering the standard ARF blind to financial fraud.

\subsection{Proposed Method: Implicit Online Oversampling}
To resolve the extreme class imbalance without introducing the unacceptable computational latency of two-stage resampling buffers, we propose a Cost-Sensitive formulation integrated directly into the ARF bagging mechanism. 

In standard Online Bagging, the weight $w$ of an instance is drawn from a Poisson distribution:
\begin{equation}
    w \sim \text{Poisson}(\lambda=1)
\end{equation}

Our proposed \textit{Implicit Online Oversampling} modifies this uniform distribution by incorporating an asymmetric cost matrix. We define the cost of a False Negative (missed fraud) as $C_{FN} = \beta$ and the cost of a False Positive (false alarm) as $C_{FP} = 1$. The parameter $\beta$ is a hyperparameter proportional to the imbalance ratio of the stream.

When a transaction $(x_t, y_t)$ arrives, the Poisson parameter $\lambda$ is dynamically adjusted based on the true class label:
\begin{equation}
    \lambda(y_t) = 
    \begin{cases} 
      1, & \text{if } y_t = 0 \text{ (Legitimate)} \\
      \beta, & \text{if } y_t = 1 \text{ (Fraud)}
    \end{cases}
\end{equation}

Thus, the instance weight is generated as:
\begin{equation}
    w \sim \text{Poisson}(\lambda(y_t))
\end{equation}

By setting $\lambda(y_t=1) = \beta$, each rare fraudulent transaction is implicitly presented to the Hoeffding Tree $\approx \beta$ times during the single-pass update. This massive injection of minority-class weight artificially amplifies the Information Gain associated with fraud-discriminating features. 

\textbf{Computational Efficiency:} Unlike traditional Synthetic Minority Over-sampling Technique (SMOTE) or batch-based oversampling, our method does not physically generate or buffer synthetic data points in memory. The oversampling is executed purely mathematically through the Poisson weight multiplier $w$ during the tree's sufficient statistics update phase. Consequently, the proposed CS-ARF maintains an $O(1)$ processing time per instance, satisfying the strict latency requirements of real-time payment networks while successfully circumventing the blindness caused by extreme class imbalance.
